\chapter{Objet et périmètre}

\section{Objet du document}

Le présent cahier décrit la stratégie de vérification de ChargeTrackr, les
environnements nécessaires, les résultats des contrôles automatisés et les
scénarios de recette manuelle. Il constitue le lien vérifiable entre le backlog,
les critères d'acceptation, le code source et le futur procès-verbal de recette.

Le document distingue volontairement :

\begin{itemize}[leftmargin=6mm]
  \item les contrôles déjà exécutés automatiquement et reproductibles ;
  \item les scénarios fonctionnels à exécuter devant l'encadrant ou le client ;
  \item les validations qui nécessitent une infrastructure de production, un
        prestataire financier ou une borne physique ;
  \item les fonctionnalités partielles ou reportées du backlog.
\end{itemize}

\section{Périmètre testé}

Le périmètre couvre les applications React et Laravel, la base PostgreSQL, le
gateway OCPP 1.6 JSON, le simulateur de bornes, l'adaptateur de paiement simulé,
les files de tâches, le temps réel et les exports documentaires. Les rôles
Visiteur, Client, Technicien, Opérateur, Administrateur et Super Administrateur
sont inclus.

\begin{table}[H]
\centering
\caption{Références du cahier de tests}
\begin{tabularx}{\textwidth}{|L{3.5cm}|Y|L{3.2cm}|}
\hline
\rowcolor{ctmint}\textbf{Référence} & \textbf{Utilisation} & \textbf{Emplacement} \\ \hline
Backlog produit & 42 User Stories et priorités & \path{docs/product-backlog.md} \\ \hline
Critères d'acceptation & Conditions vérifiables CA-xx.y & \path{docs/product-backlog-acceptance-criteria.md} \\ \hline
Suivi du backlog & État réalisé, partiel ou reporté & \path{docs/product-backlog-progress.md} \\ \hline
Contrat REST & Routes, schémas et sécurité Sanctum & \path{docs/api/openapi.json} \\ \hline
Documentation OCPP & Architecture, commandes et scénarios & \path{docs/ocpp.md} \\ \hline
Documentation paiement & Adaptateur et simulateur externe & \path{docs/payment-simulator.md} \\ \hline
\end{tabularx}
\end{table}

\section{Gestion du document}

\begin{table}[H]
\centering
\caption{Historique des versions}
\begin{tabularx}{\textwidth}{|C{1.5cm}|C{2.5cm}|L{3.4cm}|Y|}
\hline
\rowcolor{ctmint}\textbf{Version} & \textbf{Date} & \textbf{Auteur} & \textbf{Modification} \\ \hline
1.0 & 29/07/2026 & Wajdi Ben Abdeljelil & Création du cahier, inventaire des suites, résultats automatiques, recette manuelle et traçabilité complète. \\ \hline
1.1 & 17/08/2026 & Wajdi Ben Abdeljelil & Actualisation de la campagne, ajout du frontend, de l'infrastructure, du Simulation Lab, du déploiement et des scénarios de recharge simulée. \\ \hline
\end{tabularx}
\end{table}
